\section{The actual zero ideal of an original line section}
Let $X$ be any scheme, $L$ an original invertible module sheaf, and
$s:\mathcal O_X\to L$ an actual section. Retain the existing sheaf dual,
its actual contravariant map, and the original dual-unit isomorphism. Set
\[
 \epsilon_s:L^\vee\xrightarrow{s^\vee}\mathcal O_X^\vee
       \xrightarrow{\sim}\mathcal O_X,\qquad
 Z(s)=I_{\epsilon_s}.
\]
Here $I$ is the genuine quasi-coherent image ideal, not a prescribed chart ideal.
Invertibility of $L$ implies invertibility and hence quasi-coherence of the
same original $L^\vee$ by existing theorems.

\begin{lemma}[The actual zero ideal on a genuine frame]
\label{lem:line-section-zero-ideal-frame}
\leanok
\lean{EnvelopeProperty.LineSectionZeroIdeal.zeroIdeal,EnvelopeProperty.LineSectionZeroIdeal.zeroIdeal_ideal_of_frame}
\uses{thm:quasicoherent-image-ideal}
Let $X$ be any scheme, $L$ an original invertible module sheaf and
$s:\mathcal O_X\to L$ an actual section. Let $Z(s)$ be the actual image ideal
of the original dual map $s^\vee:L^\vee\to\mathcal O_X^\vee$ followed
by the existing dual-unit isomorphism. For every actual affine open $U$ and actual over-site frame
$e:L|_U\simeq\mathcal O_U$, put
$c=e(s|_U(1))\in\Gamma(X,U)$. Then the actual original affine ideal is
\[
 Z(s)(U)=(c).
\]
No integrality, nonempty-open, field or regularity assumption is imposed.
\end{lemma}
\begin{proof}
The defining original dual precomposition and dual-unit evaluation give
$\epsilon_s(\alpha)=\alpha(s|_U(1))$ for every actual functional on $U$.
The original frame $e$ itself is such a functional and its image is $c$.
The literal image is an ideal, so it contains $(c)$.
Conversely the SAME original frame inverse gives
$s|_U(1)=c\,e^{-1}(1)$. Actual linearity of any original functional
$\alpha$ gives $\epsilon_s(\alpha)=c\alpha(e^{-1}(1))$, which belongs
to $(c)$. Hence the whole original image equals $(c)$. The genuine
image-ideal affine formula identifies it with the stated actual $Z(s)(U)$.
\end{proof}

\begin{theorem}[An actual nonzero line section on an integral scheme is Cartier]
\label{thm:line-section-zero-ideal-cartier}
\leanok
\lean{EnvelopeProperty.LineSectionZeroIdeal.zeroIdeal_isInvertible_of_ne_zero}
\uses{lem:line-section-zero-ideal-frame}
Let $X$ be integral, $L$ an original invertible module sheaf and
$s:\mathcal O_X\to L$ an original nonzero section. Its original dual-map
image ideal $Z(s)$ is invertible in the existing effective-Cartier ideal
sense. No separate chart-regularity or generic-nonvanishing certificate
is required.
\end{theorem}
\begin{proof}
First the coefficient of $s$ in every frame on every nonempty open is
nonzero. Indeed, if it were zero on such an open $U$, take any other
nonempty framed open $V$. Irreducibility of the original integral scheme
makes $U\cap V$ nonempty. Injectivity of the actual frame on $U$ makes the section vector
zero there. Restrict that original vector to the intersection. Naturality
of the original section and the frame on $V$ implies that the coefficient
on $V$ restricts to zero on the same intersection. The original structure-ring restriction is injective on an
integral scheme with nonempty target open, so that coefficient is zero
on $V$. The actual line atlas covers $X$; sheaf separatedness gives the
zero original section, and the existing unit-section correspondence then
gives $s=0$, contradicting the hypothesis.
Now fix any point and refine its existing invertible-sheaf frame to an
actual affine neighbourhood $U$ of that point. This $U$ is nonempty.
Its original section ring is a domain; the derived nonzero coefficient
$c$ is therefore a nonzerodivisor. The preceding lemma identifies the
same actual ideal $Z(s)(U)$ with $(c)$. These actual affine regular
generators at every point are precisely the existing effective-Cartier
ideal definition, proving the conclusion.
\end{proof}

This constructs no canonical differential or line factory. For an original
birational morphism between smooth integral schemes of the same relative
dimension, its existing differential section is a future consumer only
after its nonzero hypothesis is genuinely derived from the original
birational map. The latter fact and its DVR discrepancy identification
are not assumptions or conclusions of the present two results.
